Imagine you are watching a stream of data go by. It might be the ups and downs of a stock market ticker, the sequence of DNA bases in a genome, or the turbulent velocity of a fluid.

At first glance, it looks random. But is it?

We are often tasked with finding \textit{patterns} in data. But what exactly is a pattern? In statistics, we talk about "correlation". In dynamical systems, we talk about "attractors". Computational Mechanics provides a different answer: a pattern is a structure that allows us to \textbf{predict}.

To understand a process is to be able to predict its future given its past. If a process is perfectly random (like a coin flip), the past tells us nothing about the future. If it is periodic (like a clock), the past tells us everything. Most interesting processes lie somewhere in between: they are \textit{structured stochastic processes}.

\begin{sidebar}{Intrinsic Computation}
Computational Mechanics views nature not just as a system of energy, but as a system of \textit{information processing}. The universe computes. A process ``stores'' information from the past (memory) and ``processes'' it to generate the future. Our goal is to reverse-engineer this computation from the data alone.
\end{sidebar}

This guide will walk you through the framework of Computational Mechanics \cite{shalizi2000computational}. We will not just learn definitions; we will construct the optimal predictive model---the $\epsilon$-machine---from scratch. We will see how to measure the "memory" of a process and verify that our models are as simple as possible, but no simpler.
